Виведення для рівняння Лапласа:

\begin{equation}
    \label{eqn:theory_thermo_laplace_eqn}
    \nabla \cdot \lambda_\theta(x) \nabla \theta(x) = 0, x \in \Omega
\end{equation}

перепишемо рівняння у вигляді:

\begin{equation}
    \lambda_\theta(x) \Delta \theta(x) = - \nabla \lambda_\theta(x) \cdot \nabla \theta(x), x \in \Omega
\end{equation}

\begin{equation}
    \Delta \theta(x) = - \frac{\nabla \lambda_\theta(x) \cdot \nabla \theta(x)}{\lambda_\theta(x)}, x \in \Omega
\end{equation}

\noindent де $\lambda_\theta(x)$ сталий на всій області $\Omega$ крім області $G$: 

\begin{equation}
    \lambda_\theta(x) = const, x \in \Omega \setminus G
\end{equation}

Частина інтегрального подання для правої частини:
\begin{equation}
    b(x) = \int \limits_{G} \frac{\nabla \lambda_\theta(\xi) \cdot \nabla \theta(\xi)}{\lambda_\theta(\xi)} 
        \hat{\theta}(x,\xi) d \xi,
\end{equation}

\noindent де $\hat{\theta}(x, \xi)$ -- фундаментальний розв'язок рівняння Лапласа (\ref{eqn:theory_thermo_laplace_eqn}). 
Продовжимо процес виведення за допомогою формули диференціювання множення для скалярної функції $p(x)$ і вектора 
$\overline{V(x)}$:
\begin{equation}
    \label{eqn:theory_vector_differentiation}
    \nabla p(x) \cdot \overline{V(x)} = \nabla \cdot (p(x)\overline{V(x)}) - p(x) \nabla \cdot \overline{V(x)}
\end{equation}

\noindent покладемо $p(x) = \theta(x)$ і 
$\overline{V(x)} = \nabla \lambda_{\theta}(x) \frac{\hat{\theta}(x, \xi)}{\lambda_{\theta}(x)}$. Отримаємо:

\begin{equation}
    \begin{split}
        b(x) =  \int \limits_{G} \nabla \cdot 
            \left( \theta(\xi) \frac{\nabla \lambda_\theta(\xi) \hat{\theta}(x, \xi)}{\lambda_\theta(\xi)} \right) d \xi 
            -
            \int \limits_{G} \theta(\xi) \nabla \cdot  
                \left( \frac{\nabla \lambda_\theta(\xi) \hat{\theta}(x, \xi)}{\lambda_\theta(\xi)} \right) d \xi 
    \end{split}
\end{equation}

\noindent До першого доданка застосуємо формулу Ґріна, а у другому виконаємо диференціювання:

\begin{equation}
    \begin{split}
        b(x) =  \int \limits_{\Gamma_G} \overline{n} 
            \left[ \theta(\xi) \frac{\nabla \lambda_\theta(\xi) \hat{\theta}(x, \xi)}{\lambda_\theta(\xi)} 
                \right] d \xi - \\
            -
            \int \limits_{G} \theta(\xi) \left[
                \Delta \lambda_{\theta}(\xi) \frac{\hat{\theta}(x, \xi)}{\lambda_{\theta}(\xi)} +
                \nabla \left(\frac{\hat{\theta}(x, \xi)}{\lambda_{\theta}(\xi)}\right) 
                    \cdot \nabla \lambda_{\theta}(\xi) \right] d \xi 
    \end{split}
\end{equation}

\noindent $\Gamma_G$ -- межа області G. Перший доданок рівний нулю оскільки містить похідну від сталої функції на межі 
$\Gamma_G$. Виконаємо диференціювання у другому доданку і отримаємо:

\begin{equation}
    \begin{split}
        b(x) =  - \int \limits_{G} \theta(\xi) \frac{\Delta \lambda_{\theta}(\xi)}{\lambda_{\theta}(\xi)} 
            \hat{\theta}(x, \xi) d \xi +
            \int \limits_{G} \theta(\xi) \frac{\nabla \lambda_{\theta}(\xi) \cdot \nabla \lambda_{\theta}(\xi)}
                {\lambda_{\theta}^2(\xi)}\hat{\theta}(x, \xi) d \xi - \\
            -  
            \int \limits_{G} \theta(\xi) \frac{\nabla \lambda_{\theta}(\xi) \cdot \nabla \hat{\theta}(x, \xi)}
                {\lambda_{\theta}(\xi)} d \xi
    \end{split}
\end{equation}

\noindent Розіб'ємо $G$ на $m$ трикутних елементів $G_i$, 
$G = \bigcup_{i=1}^m G_i$, $G_i \cap G_j = \varnothing, \forall i \neq j$. На кожному із цих елементів апроксимуємо 
$\theta(x)$ сталим значенням у центрі мас трикутника $x_i$, $\theta_i = \theta(x_i)$. Перепишемо попереднє рівняння із
врахуванням апроксимацій:

\begin{equation}
    \begin{split}
        b(x) =  - \sum_{j=1}^{m} \theta_j \int \limits_{G_j} \frac{\Delta \lambda_{\theta}(\xi)}{\lambda_{\theta}(\xi)} 
            \hat{\theta}(x, \xi) d \xi +
            \sum_{j=1}^{m} \theta_j \int \limits_{G_j} \frac{\nabla \lambda_{\theta}(\xi) \cdot \nabla 
                \lambda_{\theta}(\xi)} {\lambda_{\theta}^2(\xi)}\hat{\theta}(x, \xi) d \xi - \\
            -  
            \sum_{j=1}^{m} \theta_j \int \limits_{G_j} \frac{\nabla \lambda_{\theta}(\xi) \cdot \nabla 
                \hat{\theta}(x, \xi)} {\lambda_{\theta}(\xi)} d \xi
    \end{split}
\end{equation}

\noindent До третього доданка застосуємо диференціальну формулу (\ref{eqn:theory_vector_differentiation}). Покладемо
$p(x) = \hat{\theta}(x, \xi)$, $\overline{V} = \frac{\nabla \lambda_{\theta}(x)}{\lambda_{\theta}(x)}$:

\begin{equation}
    \begin{split}
        \int \limits_{G_j} \frac{\nabla \lambda_{\theta}(\xi) \cdot \nabla \hat{\theta}(x, \xi)} 
            {\lambda_{\theta}(\xi)} d \xi = 
        \int \limits_{G_j} \nabla \cdot \left[\frac{\hat{\theta}(x, \xi) \nabla \lambda_{\theta}(\xi)}
            {\lambda_{\theta}(\xi)} \right] d \xi - \\ 
        - 
        \int \limits_{G_j} \hat{\theta}(x, \xi) \nabla \cdot \left[ \frac{\nabla \lambda_{\theta}(\xi)}
            {\lambda_{\theta}(\xi)} \right] d \xi
    \end{split}
\end{equation}

\noindent Застосуємо формулу Ґріна до першого доданка і інтегруємо другий доданок:

\begin{equation}
    \begin{split}
        \int \limits_{G_j} \frac{\nabla \lambda_{\theta}(\xi) \cdot \nabla \hat{\theta}(x, \xi)} 
            {\lambda_{\theta}(\xi)} d \xi = 
        \int \limits_{\Gamma_{G_j}} \overline{n} \frac{\nabla \lambda_{\theta}(\xi)}
            {\lambda_{\theta}(\xi)} \hat{\theta}(x, \xi) ds - \\ 
        - 
        \int \limits_{G_j} \frac{\Delta \lambda_{\theta}(\xi)}{\lambda_{\theta}(\xi)} \hat{\theta}(x, \xi)  d \xi
        + 
        \int \limits_{G_j} \frac{\nabla \lambda_{\theta}(\xi) \cdot \nabla \lambda_{\theta}(\xi)}
            {\lambda_{\theta}^2(\xi)} \hat{\theta}(x, \xi) d \xi
        \end{split}
\end{equation}

\noindent Підставимо все у формулу для правої частини:

\begin{equation}
    \begin{split}
        b(x) =  - \sum_{j=1}^{m} \theta_j \int \limits_{\Gamma_{G_j}} \overline{n} \frac{\nabla \lambda_{\theta}(\xi)}
        {\lambda_{\theta}(\xi)} \hat{\theta}(x, \xi) ds
    \end{split}
\end{equation}

